\documentclass{article}

\usepackage[spanish]{babel}
\usepackage{amsmath,amssymb}

\newcommand{\N}{\mathbb{N}}
\newcommand{\Z}{\mathbb{Z}}
\newcommand{\Q}{\mathbb{Q}}
\newcommand{\R}{\mathbb{R}}
\newcommand{\C}{\mathbb{C}}

\title{Certamen 4 - Cáculo II}
\author{Benjamín Rivas Beltrán}
\date{\today}

\begin{document}

\maketitle

\section{Descripción de conjunto $D$}
\[
  D = \{(x, y) \in \R^2 : 4 < (x - 1)^2 + y^2 \leq 9 \}
\]

\begin{itemize}
  \item Puntos interiores:
  \[
    D^\circ = \{(x, y) \in \R^2 : 4 < (x - 1)^2 + y^2 < 9 \}
  \]
  
  \item Puntos adherentes:
  \[
    \overline{D} = \{(x, y) \in \R^2 : 4 \leq (x - 1)^2 + y^2 \leq 9 \}
  \]
  
  \item Puntos frontera:
  \[
    \overline{D} = \{(x, y) \in \R^2 : (x - 1)^2 + y^2 = 4 \} \ \cup \ \{(x, y) \in \R^2 : (x - 1)^2 + y^2 = 9 \} \ \cup \ {(1, 0)}
  \]

  \item Abierto:
  \[
    D^\circ \neq D \Rightarrow D \text{ no es abierto}
  \]

  \item Cerrado:
  \[
    \overline{D} \neq D \Rightarrow D \text{ no es cerrado}
  \]

  \item Compacto:
  \[
    D \text{ no es cerrado} \Rightarrow D \text{ no es compacto}
  \]
\end{itemize}

\section{Análisis de función}

\[
f(x, y) = \sqrt{\frac{x^2 + y^2 + 1}{x - y}}
\]

\begin{itemize}
  \item Dominio:
    \begin{align*}
      Dom(f) &= \{(x, y) \in \R^2 : f(x, y) \in \R \} \\
      &= \{(x, y) \in \R^2 : \sqrt{\frac{x^2 + y^2 + 1}{x - y}} \in \R \} \\
      &= \{(x, y) \in \R^2 : \frac{x^2 + y^2 + 1}{x - y} \geq 0 \} \\
      &= \{(x, y) \in \R^2 : x - y \geq 0 \} \\
      &= \{(x, y) \in \R^2 : x \geq y \} \\
    \end{align*}
    
  \item Curvas de nivel:
    \[
      C_k = \{(x, y) \in Dom(f) : f(x, y) = k\}
    \]

    \begin{itemize}
      \item \begin{align*}
        C_{-2} &= \{(x, y) \in Dom(f) : f(x, y) = -2\} \\
        &= \{(x, y) \in Dom(f) : \sqrt{\frac{x^2 + y^2 + 1}{x - y}} = -2\} \\
        &= \emptyset
      \end{align*}

      La curva de nivel es el conjunto vacío, pues una raíz cuadrada no puede ser negativa.

      \item \begin{align*}
        C_{\sqrt{2}} &= \{(x, y) \in Dom(f) : f(x, y) = \sqrt{2}\} \\
        &= \{(x, y) \in Dom(f) : \sqrt{\frac{x^2 + y^2 + 1}{x - y}} = \sqrt{2}\} \\
        &= \{(x, y) \in Dom(f) : \frac{x^2 + y^2 + 1}{x - y} = 2\} \\
        &= \{(x, y) \in Dom(f) : x^2 + y^2 + 1 = 2x - 2y\} \\
        &= \{(x, y) \in Dom(f) : x^2 - 2x + y^2 + 2y + 1 = 0\} \\
        &= \{(x, y) \in Dom(f) : x^2 - 2x + 1 + y^2 + 2y + 1 = 1\} \\
        &= \{(x, y) \in Dom(f) : (x - 1)^2 + (y + 1)^2 = 1\} \\
      \end{align*}

      La curva de nivel es un circunferencia de radio 1 centrada en el punto $(1, -1)$.

      \item \begin{align*}
        C_{2} &= \{(x, y) \in Dom(f) : f(x, y) = 2\} \\
        &= \{(x, y) \in Dom(f) : \sqrt{\frac{x^2 + y^2 + 1}{x - y}} = 2\} \\
        &= \{(x, y) \in Dom(f) : \frac{x^2 + y^2 + 1}{x - y} = 4\} \\
        &= \{(x, y) \in Dom(f) : x^2 + y^2 + 1 = 4x - 4y\} \\
        &= \{(x, y) \in Dom(f) : x^2 - 4x + y^2 + 4y + 1 = 0\} \\
        &= \{(x, y) \in Dom(f) : x^2 - 4x + 4 + y^2 + 2y + 4 = 7\} \\
        &= \{(x, y) \in Dom(f) : (x - 2)^2 + (y + 2)^2 = 7\} \\
      \end{align*}

      La curva de nivel es un circunferencia de radio $\sqrt{7}$ centrada en el punto $(2, -2)$.

    \end{itemize}
\end{itemize}

\section{Continuidad y derivadas direccionales}

\[
  f(x, y) = \begin{cases}
    \frac{5x^2y}{x^2 + y^2}, & \text{si } (x, y) \neq (0, 0) \\
    0, & \text{si } (x, y) = (0, 0) 
  \end{cases}
\]

\begin{itemize}
  \item Continuidad
  
  \begin{itemize}
    \item Caso 1 ($(x, y) \neq (0, 0)$):
      Dado que para todo $(x, y) \in \R^2 : (x, y) \neq (0, 0)$, $x^2 + y^2 \neq 0$, por álgebra de funciones continuas, $f$ es continua en todos esos puntos.

    \item Caso 2 ($(x, y) = (0, 0)$):
      Se sabe que $f(0, 0) = 0$. Entonces, para que $f$ sea continua en ese punto se debe cumplir que:
      \[
        \lim_{(x, y) \to (0, 0)} f(x, y) = 0
      \]

      Desarrollando:
      \[
        \lim_{(x, y) \to (0, 0)} f(x, y) = \lim_{(x, y) \to (0, 0)} \frac{5x^2y}{x^2 + y^2}
      \]

      Note que:
      \begin{align*}
      0 \leq \left| \frac{5x^2y}{x^2 + y^2} \right| &= \frac{5x^2 \left| y \right|}{\sqrt{x^2 + y^2} \sqrt{x^2 + y^2}} \\ &= \frac{5x^2}{\sqrt{x^2 + y^2} }\frac{\left| y \right|}{\sqrt{x^2 + y^2}} \\ 
      &\leq \frac{5x^2}{\sqrt{x^2} }\frac{\left| y \right|}{\sqrt{y^2}} \\
      &= \frac{5x^2}{\left|x\right|}\frac{\left| y \right|}{\left|y\right|} \\
      &= 5\left|x\right|
      \end{align*}

      Y como:
      \[
        \lim_{(x, y) \to (0, 0)} 0 = \lim_{(x, y) \to (0, 0)} 5 \left|x\right| = 0
      \]

      Por el teorema del sandwich, se concluye que:
      \[
      \lim_{(x, y) \to (0, 0)} \frac{5x^2y}{x^2 + y^2} = 0
      \]

      Luego, $f$ es continua en $(0, 0)$.
  \end{itemize}

  Dado que $f$ es continua en todo su dominio, se concluye que es una función continua
  
  
  \item Derivadas parciales en $0, 0$
    \begin{align*}
      \frac{\partial f}{\partial x}(0, 0) &= \lim_{h \to 0} \frac{f((0, 0) + (h, 0)) - f(0, 0)}{h} \\
      &= \lim_{h \to 0} \frac{f(h, 0) - 0}{h} \\
      &= \lim_{h \to 0} \frac{\frac{5h^2 0}{h^2 + 0^2} - 0}{h} \\
      &= \lim_{h \to 0} \frac{0}{h} \\
      &= 0
    \end{align*}

    \begin{align*}
      \frac{\partial f}{\partial y}(0, 0) &= \lim_{h \to 0} \frac{f((0, 0) + (0, h)) - f(0, 0)}{h} \\
      &= \lim_{h \to 0} \frac{f(0, h) - 0}{h} \\
      &= \lim_{h \to 0} \frac{\frac{5(0)^2 h}{0^2 + y^2} - 0}{h} \\
      &= \lim_{h \to 0} \frac{0}{h} \\
      &= 0
    \end{align*}
\end{itemize}

\section{Puntos críticos}

\[
f(x, y) = x^3 - y^3 - 12x + 27y + 35
\]

Calculando las derivadas parciales:
\[
f_x(x, y) = 3x^2 - 12
\]
\[
f_y(x, y) = -3y^2 + 27
\]

Puntos críticos:
\begin{equation}
  3x^2 - 12 = 0
\end{equation}

\begin{equation}
  -3y^2 + 27 = 0
\end{equation} 

Resolviendo:
\begin{align*}
  3x^2 - 12 = 0 \iff x = 2 \vee x = -2 \\
  -3y^2 + 27 = 0 \iff y = 3 \vee x = -3
\end{align*}

Los puntos críticos son:
\begin{align*}
  P_1(2, 3) \\
  P_2(2, -3) \\
  P_3(-2, 3) \\
  P_4(-2, -3)
\end{align*}

Calculando las segundas derivadas:
\begin{align*}
f_{xx}(x, y) = 6x \\
f_{yy}(x, y) = -6y \\
f_{xy}(x, y) = 0 \\
f_{yx}(x, y) = 0
\end{align*}

El determinante de la matriz Hessiana es:
\[
D(x, y) = f_{xx}(x, y)f_{yy}(x, y) - f_{xy}(x, y)f_{yx}(x, y) = 6x(-6y) - 0 = -36xy
\]

Resolviendo para cada punto crítico:
\begin{itemize}
  \item $P_1(2, 3)$:
  \[
    D(2, 3) = -36(2)(3) = -216
  \]
  \item $P_2(2, -3)$:
  \[
    D(2, -3) = -36(2)(-3) = 216
  \]
  \[
    f_{xx}(2, -3) = 6(2) = 12 > 0 \Rightarrow P_2(2, -3) \text{ es un mínimo local}
  \]
  \item $P_3(-2, 3)$:
  \[
    D(-2, 3) = -36(-2)(3) = 216
  \]
  \[
    f_{xx}(-2, 3) = 6(-2) = -12 < 0 \Rightarrow P_3(-2, 3) \text{ es un máximo local}
  \]
  \item $P_4(-2, -3)$:
  \[
    D(-2, -3) = -36(-2)(-3) = -216
  \]
\end{itemize}

Por lo tanto, los puntos críticos son:
\begin{itemize}
  \item $P_1(2, 3)$: punto silla
  \item $P_2(2, -3)$: máximo local
  \item $P_3(-2, 3)$: mínimo local
  \item $P_4(-2, -3)$: punto silla
\end{itemize}


\end{document}